\documentclass[aspectratio=169,11pt]{beamer}
\usepackage[bahasa]{babel}
\usepackage[utf8]{inputenc}
\usepackage[T1]{fontenc}
\usepackage{tabularx}
\usepackage{booktabs}
\usepackage{tikz}
\usepackage{graphicx}
\usepackage[most]{tcolorbox}
\usepackage{amsmath}
\usepackage{amssymb}
\usetikzlibrary{positioning, arrows.meta, shapes.geometric, backgrounds, fit, calc}

% Warna Korporat PLN + Premium Accents
\definecolor{plnblue}{RGB}{0,75,135}
\definecolor{plnorange}{RGB}{255,107,53}
\definecolor{plncyan}{RGB}{0,163,224}
\definecolor{plngreen}{RGB}{0,150,136}
\definecolor{lightgrey}{RGB}{248,249,250}
\definecolor{darkgrey}{RGB}{50,50,50}

% TikZ BPMN Styles
\tikzset{
    startstyle/.style={circle, draw=red!70!black, fill=red!15, minimum size=0.9cm, align=center, font=\tiny\bfseries, line width=1pt},
    endstyle/.style={circle, draw=red!70!black, fill=red!25, minimum size=0.9cm, align=center, font=\tiny\bfseries, line width=1.5pt},
    startstyle_tobe/.style={circle, draw=green!60!black, fill=green!15, minimum size=0.9cm, align=center, font=\tiny\bfseries, line width=1pt},
    endstyle_tobe/.style={circle, draw=plnblue!80!black, fill=plnblue!25, minimum size=0.9cm, align=center, font=\tiny\bfseries, line width=1.5pt},
    manual/.style={rectangle, rounded corners=2pt, draw=gray!70, fill=gray!8, text width=2.4cm, minimum height=1.2cm, align=center, font=\tiny, line width=0.6pt},
    waiting/.style={rectangle, rounded corners=2pt, draw=red!50, fill=red!5, text width=2.4cm, minimum height=1.2cm, align=center, font=\tiny\itshape, line width=0.6pt},
    autostep/.style={rectangle, rounded corners=2pt, draw=plnblue!80, fill=plnblue!12, text width=2.4cm, minimum height=1.2cm, align=center, font=\tiny\bfseries, line width=0.6pt},
    humanstep/.style={rectangle, rounded corners=2pt, draw=plnorange, fill=plnorange!10, text width=2.4cm, minimum height=1.2cm, align=center, font=\tiny, line width=0.6pt},
    gateway/.style={diamond, draw=gray!70, fill=yellow!8, minimum width=1.1cm, minimum height=1.1cm, align=center, font=\tiny\bfseries, inner sep=1pt, aspect=1},
    gateway_tobe/.style={diamond, draw=plnorange!80, fill=plnorange!8, minimum width=1.1cm, minimum height=1.1cm, align=center, font=\tiny\bfseries, inner sep=1pt, aspect=1},
    arr/.style={-{Stealth[length=5pt]}, thick, color=gray!60},
    arr_tobe/.style={-{Stealth[length=5pt]}, thick, color=plnblue!70},
    arrok/.style={-{Stealth[length=5pt]}, thick, color=green!60!black},
    arrno/.style={-{Stealth[length=5pt]}, thick, color=red!50, dashed},
    timelbl/.style={font=\tiny\bfseries, text=red!60!black},
    timelbl_tobe/.style={font=\tiny\bfseries, text=green!50!black},
    lanelab/.style={font=\tiny\bfseries, text=gray!70!black, rotate=90}
}


% Konfigurasi Tema Beamer Premium
\mode<presentation>{
    \usetheme{Madrid}
    \usecolortheme[named=plnblue]{structure}
    \setbeamercolor{palette primary}{bg=plnblue,fg=white}
    \setbeamercolor{palette secondary}{bg=plnorange,fg=white}
    \setbeamercolor{title}{bg=plnblue,fg=white}
    \setbeamercolor{frametitle}{bg=plnblue,fg=white}
    \setbeamerfont{frametitle}{series=\bfseries, size=\large}
    \setbeamercolor{background canvas}{bg=white}
    \setbeamertemplate{navigation symbols}{}
    \setbeamertemplate{items}[circle]
}

% Custom Boxes
\newtcolorbox{premiumbox}[1][]{%
  colback=plnblue!5!white, colframe=plnblue, fonttitle=\bfseries\small,
  sharp corners, boxrule=1pt, left=4pt, right=4pt, top=4pt, bottom=4pt,
  title={#1}, before skip=4pt, after skip=10pt, breakable
}

\newtcolorbox{orangebox}[1][]{%
  colback=plnorange!5!white, colframe=plnorange, fonttitle=\bfseries\small,
  sharp corners, boxrule=1pt, left=4pt, right=4pt, top=4pt, bottom=4pt,
  title={#1}, before skip=4pt, after skip=10pt, breakable
}

\newtcolorbox{greenbox}[1][]{%
  colback=plngreen!5!white, colframe=plngreen, fonttitle=\bfseries\small,
  sharp corners, boxrule=1pt, left=4pt, right=4pt, top=4pt, bottom=4pt,
  title={#1}, before skip=4pt, after skip=10pt, breakable
}

\title{\LARGE Optimalisasi Prosedur \& Alur Kerja Terintegrasi}
\subtitle{Modul Pelatihan Level 4 (Mastery) -- Transformasi Tata Kelola Kearsipan\\ \textbf{PT PLN (Persero)}}
\author{Fasilitator Utama}
\institute[PLN]{Center of Excellence Records Management}
\date{April 2026}

% Makro untuk Icon Sederhana
\newcommand{\stepicon}[2]{
    \begin{tikzpicture}
        \node[circle, fill=#1, text=white, minimum size=0.6cm, inner sep=0, font=\tiny\bfseries] {#2};
    \end{tikzpicture}
}

\begin{document}

% SLIDE 1: TITLE PAGE
{
\beamertemplatenavigationsymbolsempty
\begin{frame}[plain]
    \titlepage
\end{frame}
}

% SLIDE 2: VISION & CONCEPT
\begin{frame}{Visi Masa Depan: Ekosistem Kearsipan Digital PLN}
    \begin{columns}
        \column{0.55\textwidth}
        \begin{premiumbox}[Transformasi Budaya \& Teknologi]
            \small
            \textit{"Archiving is no longer a post-action task, but an embedded engine of business reliability."}
            \vspace{0.3cm}
            \begin{itemize}
                \item \textbf{Embedded:} Menyatu dalam transaksi operasional.
                \item \textbf{Autonomous:} Tarik data (\textit{auto-capture}) tanpa manual.
                \item \textbf{Authentic:} Menjamin bukti legal yang tak terbantahkan.
            \end{itemize}
        \end{premiumbox}
        \centering
        \stepicon{plnblue}{1} Diagnosa \quad \stepicon{plncyan}{2} Desain \quad \stepicon{plnorange}{3} Deliver
        
        \column{0.45\textwidth}
        \centering
        \includegraphics[width=\textwidth]{/home/lambda_one/.gemini/antigravity/brain/ef68cde3-b0c3-4a08-91b5-b3462ac425ba/pln_digital_archive_concept_1776762781196.png}
    \end{columns}
\end{frame}

% SLIDE 3: PROBLEM QUANTIFICATION
\begin{frame}{Biaya Tersembunyi dari Prosedur Manual (As-Is)}
    \begin{columns}
        \column{0.5\textwidth}
        \begin{orangebox}[Analisis Inefisiensi SPK/Kontrak]
            \small
            \begin{description}
                \item[\textcolor{red}{Waktu Tunggu}] 3--5 Hari Kerja.
                \item[\textcolor{red}{Human Error}] 35\% Metadata Salah Input.
                \item[\textcolor{red}{Risiko Audit}] Dokumen tercecer/sulit dicari.
                \item[\textcolor{red}{Biaya Fisik}] Kertas, Kurir, \& Ruang Simpan.
            \end{description}
        \end{orangebox}
        
        \column{0.5\textwidth}
        \centering
        \textbf{Kerugian Simulasi (UIT GI 150kV)}
        \vspace{0.3cm}
        \begin{tikzpicture}
            \draw[lightgrey, fill=lightgrey] (0,0) rectangle (4,2.5);
            \node at (2,1.8) {\small \textbf{Potensi Kerugian}};
            \node at (2,1.2) {\huge \textcolor{red}{\textbf{Rp 1,04 M}}};
            \node at (2,0.5) {\tiny \textbf{Per Tahun / Unit Kerja}};
            \draw[red, line width=1pt] (0.5,0.8) -- (3.5,0.8);
        \end{tikzpicture}
    \end{columns}
\end{frame}

% SLIDE 4: THE GAP (Friction Points)
\begin{frame}{Mengapa Proses Kita Terasa Lambat?}
    \centering
    \begin{tikzpicture}[node distance=1.2cm]
        \node[draw, fill=plnblue!10, text width=2.5cm, align=center, font=\tiny] (s1) {Penciptaan\\(Sistem A)};
        \node[draw, red, fill=red!10, right=of s1, text width=2.5cm, align=center, font=\tiny] (f1) {CETAK FISIK\\(Friction 1)};
        \node[draw, red, fill=red!10, right=of f1, text width=2.5cm, align=center, font=\tiny] (f2) {TTD MANUAL\\(Friction 2)};
        \node[draw, red, fill=red!10, right=of f2, text width=2.5cm, align=center, font=\tiny] (f3) {SCAN \& UPLOAD\\(Friction 3)};
        \node[draw, fill=plnblue!10, right=of f3, text width=2.5cm, align=center, font=\tiny] (s2) {Penyimpanan\\(Sistem B)};
        \draw[->, thick] (s1) -- (f1); \draw[->, thick] (f1) -- (f2); \draw[->, thick] (f2) -- (f3); \draw[->, thick] (f3) -- (s2);
    \end{tikzpicture}
    \vspace{0.6cm}
    \begin{premiumbox}[Diagnosis]
        \small Masalah utama pada \textbf{Gap Prosedur} yang mewajibkan konversi digital-analog-digital secara berulang.
    \end{premiumbox}
\end{frame}

% SLIDE 5: METHODOLOGY (PARBICA)
\begin{frame}{Metodologi: Standar Internasional PARBICA}
    \begin{columns}
        \column{0.5\textwidth}
        \begin{premiumbox}[PARBICA Toolkit Guideline 13]
            \small
            Acuan resmi \textit{International Council on Archives} (ICA) untuk organisasi dalam \textbf{Masa Transisi} menuju digital.
            \begin{itemize}
                \item Fokus pada Kesiapan (\textit{Readiness}).
                \item Diagnosis 7 Domain Kunci.
            \end{itemize}
        \end{premiumbox}
        
        \column{0.5\textwidth}
        \textbf{7 Domain Evaluasi (Skala 1--5)}
        \vspace{0.2cm}
        \tiny
        \begin{enumerate}
            \item \textbf{Kebijakan:} Dasar hukum di titik transaksi.
            \item \textbf{Tata Kelola:} Otoritas \& RACI Chart.
            \item \textbf{SDM:} Kapasitas \& Budaya kerja.
            \item \textbf{Interoperabilitas:} Komunikasi antar sistem.
            \item \textbf{Proses Bisnis:} Alur tanpa hambatan.
            \item \textbf{Keamanan:} Integritas data.
            \item \textbf{Preservasi:} Kelangsungan jangka panjang.
        \end{enumerate}
    \end{columns}
\end{frame}

% SLIDE 6: THEORY - CONTINUUM MODEL
\begin{frame}{Landasan Teori Kearsipan: \textit{Records Continuum Model}}
    \begin{columns}
        \column{0.55\textwidth}
        \begin{premiumbox}[Pergeseran Paradigma (Frank Upward, 1996)]
            \small
            Meninggalkan \textbf{Life-Cycle Theory} yang menganggap pelestarian sebagai tahap akhir, menuju pendekatan \textbf{Continuum}.
            \vspace{0.2cm}
            \begin{itemize}
                \item \textbf{Create \& Capture:} Penciptaan dan penangkapan arsip terintegrasi secara proaktif saat transaksi \textit{real-time}.
                \item \textbf{Organize:} Integritas dan interoperabilitas sistem terjamin sejak rancangan awal (\textit{by design}).
                \item \textbf{Pluralize:} Bukti akuntabilitas institusional yang dapat diakses lintas batas secara \textit{scalable}.
            \end{itemize}
        \end{premiumbox}
        
        \column{0.45\textwidth}
        \begin{orangebox}[Relevansi Kritis bagi PLN]
            \small Tata kelola asrip PLN tidak lagi beroperasi secara pasif di "gudang belakang" (\textit{back-end}), melainkan ditanamkan langsung (\textit{embedded}) layaknya sebuah organ pada aplikasi operasional lini depan (\textit{front-end}).
        \end{orangebox}
    \end{columns}
\end{frame}

% SLIDE 7: THEORY - GOVERNANCE & ISO
\begin{frame}{Landasan Tata Kelola Hukum \& Internasional: ISO 15489}
    \begin{columns}
        \column{0.6\textwidth}
        \begin{premiumbox}[\textit{Information Governance} (Tata Kelola Informasi)]
            \small
            \textit{"Information is a formalized asset that requires executive accountability."} (ARMA International)
            \vspace{0.2cm}
            \begin{itemize}
                \item Memastikan data tidak hanya sekadar "Tersimpan", tapi juga terjamin \textbf{Keabsahan}, \textbf{Perlindungan Hak Akses}, dan \textbf{Nilai Guna Audit}.
            \end{itemize}
        \end{premiumbox}
        \vspace{0.2cm}
        \begin{greenbox}[ISO 15489-1:2016 (Records Management)]
            \small
            \textbf{4 Karakteristik Konstruksi Arsip Digital:}
            \vspace{0.1cm}
            \begin{description}
                \tiny
                \item[Authenticity] Terbukti keaslian pencipta \& otoritas rilis.
                \item[Reliability] Akurat merepresentasikan peristiwa transaksi riil.
                \item[Integrity] Bebas dari alterasi (\textit{Hash-lock} Anti-Tamper).
                \item[Usability] Terstruktur \& mudah didiscoveri.
            \end{description}
        \end{greenbox}
        
        \column{0.4\textwidth}
        \centering
        \begin{tikzpicture}
            \node[circle, fill=plnblue!10, draw=plnblue, text width=2cm, align=center, font=\tiny\bfseries] (c1) {Accountability};
            \node[circle, fill=plnorange!10, draw=plnorange, text width=2cm, align=center, font=\tiny\bfseries, below left=0.5cm and -0.8cm of c1] (c2) {Integrity};
            \node[circle, fill=plngreen!10, draw=plngreen, text width=2cm, align=center, font=\tiny\bfseries, below right=0.5cm and -0.8cm of c1] (c3) {Compliance};
            \draw[thick, gray, -] (c1) -- (c2) -- (c3) -- (c1);
            \node[font=\bfseries, text=plnblue, align=center] at ($(c1)!0.5!(c2)!0.5!(c3)$) {\tiny TRUST};
        \end{tikzpicture}
    \end{columns}
\end{frame}

% SLIDE 8: STRATEGIC PILLARS
\begin{frame}{3 Pilar Strategis Optimalisasi Prosedur}
    \vspace{0.4cm}
    \begin{columns}[T]
        \column{0.31\textwidth}
        \centering
        \stepicon{plnblue}{1} \\ \textbf{Value-Added Focus} \\
        \vspace{0.2cm}
        \tiny Mengeliminasi aktivitas non-produktif (seperti bongkar gudang fisik).
        
        \column{0.31\textwidth}
        \centering
        \stepicon{plncyan}{2} \\ \textbf{Single Source of Truth} \\
        \vspace{0.2cm}
        \tiny Satu kali input di sistem operasional, mengalir otomatis ke arsip.
        
        \column{0.31\textwidth}
        \centering
        \stepicon{plnorange}{3} \\ \textbf{Rule-Based Routing} \\
        \vspace{0.2cm}
        \tiny Distribusi persetujuan otomatis berdasarkan kategori risiko.
    \end{columns}
\end{frame}

% SLIDE 7: INVESTIGATIVE PHASE
\begin{frame}{Proses Diagnosa: Mencari Fakta Lapangan}
    \vspace{0.3cm}
    \begin{columns}
        \column{0.33\textwidth}
        \centering
        \stepicon{darkgrey}{A} \\ \textbf{Document Shadowing} \\
        \tiny Membuntuti satu dokumen fisik hingga tiba di repositori pusat.
        
        \column{0.33\textwidth}
        \centering
        \stepicon{darkgrey}{B} \\ \textbf{Log Analysis} \\
        \tiny Membandingkan \textbf{Durasi SLA} vs \textbf{Fakta Riil} jam persetujuan manajer.
        
        \column{0.33\textwidth}
        \centering
        \stepicon{darkgrey}{C} \\ \textbf{Gemba Walk} \\
        \tiny Bertanya langsung kepada personel operasional mengenai hambatan nyata.
    \end{columns}
\end{frame}

% SLIDE 8: AS-IS BPMN (FULL)
\begin{frame}{Model BPMN Kondisi Saat Ini (As-Is)}
    \centering
    \resizebox{0.95\textwidth}{!}{%
    \begin{tikzpicture}[node distance=1.2cm and 0.8cm]
    % Lane backgrounds
    \fill[gray!5] (0.2, 2.5) rectangle (18, 0.6);
    \fill[yellow!4] (0.2, 0.4) rectangle (18, -2.3);
    \fill[gray!5] (0.2, -2.5) rectangle (18, -5.4);
    % Lane Dividers
    \draw[gray!30, thick] (0.2, 0.5) -- (18, 0.5);
    \draw[gray!30, thick] (0.2, -2.4) -- (18, -2.4);
    % Lane labels
    \node[lanelab] at (0.4, 1.6) {Staf Cabang};
    \node[lanelab] at (0.4, -1.0) {Manajer UPT};
    \node[lanelab] at (0.4, -4.0) {Kearsipan};

    \node[startstyle] at (1.5, 1.6) (s1) {Start};
    \node[manual, right=of s1] (m1) {Cetak\\Kontrak\\Rangkap 3};
    \node[manual, right=of m1] (m2) {Isi Form\\Pengantar};
    \node[manual, right=of m2] (m3) {Kirim Kurir\\Internal};
    
    \node[waiting] at (8.5, -1.0) (w1) {Antrian\\Meja};
    \node[manual, right=of w1] (m4) {Review\\\& TTD\\Basah};
    \node[gateway, right=of m4] (g1) {$\times$};
    \node[manual, right=of g1] (m5) {Stempel\\Basah};
    
    \node[manual] at (12.5, -4.0) (m6) {Terima\\Fisik};
    \node[manual, right=of m6] (m7) {Catat Buku\\Induk};
    \node[endstyle, right=of m7] (e1) {End};

    \draw[arr] (s1) -- (m1); \draw[arr] (m1) -- (m2); \draw[arr] (m2) -- (m3);
    \draw[arr] (m3.south) -- ++(0,-0.4) -| (w1.north);
    \draw[arr] (w1) -- (m4); \draw[arr] (m4) -- (g1);
    \draw[arrno] (g1.north) -- ++(0,0.8) -| node[above, font=\tiny\itshape, text=red] {Revisi} (m2.north);
    \draw[arr] (g1) -- node[above, font=\tiny\itshape] {OK} (m5);
    \draw[arr] (m5.south) -- ++(0,-0.4) -| (m6.north);
    \draw[arr] (m6) -- (m7); \draw[arr] (m7) -- (e1);

    \node[timelbl, below=0.15cm of w1] {1--3 hari!};
    \node[font=\small\bfseries, text=red!70!black] at (9, -6.5) {SIKLUS MANUAL: 3--5 HARI KERJA};
    \end{tikzpicture}
    }
\end{frame}

% SLIDE 9: TO-BE BPMN (FULL)
\begin{frame}{Model BPMN Kondisi Target (To-Be)}
    \centering
    \resizebox{0.95\textwidth}{!}{%
    \begin{tikzpicture}[node distance=1.2cm and 0.8cm]
    % Lane backgrounds
    \fill[plnblue!4] (0.2, 2.5) rectangle (18, 0.6);
    \fill[plnorange!4] (0.2, 0.4) rectangle (18, -2.3);
    \fill[plnblue!4] (0.2, -2.5) rectangle (18, -5.4);
    % Lane Dividers
    \draw[plnblue!20, thick] (0.2, 0.5) -- (18, 0.5);
    \draw[plnblue!20, thick] (0.2, -2.4) -- (18, -2.4);
    % Lane labels
    \node[lanelab, text=plnblue!80!black] at (0.4, 1.6) {ERP / SAP};
    \node[lanelab, text=plnblue!80!black] at (0.4, -1.0) {Pimpinan};
    \node[lanelab, text=plnblue!80!black] at (0.4, -4.0) {Repositori};

    \node[startstyle_tobe] at (1.5, 1.6) (s2) {Event};
    \node[autostep, right=of s2] (a1) {Auto-Capture\\Metadata\\ERP};
    \node[autostep, right=of a1] (a2) {Generate\\Draft Digital\\(PDF/A)};
    
    \node[humanstep] at (7.5, -1.0) (h1) {Digital Review\\via E-Office};
    \node[gateway_tobe, right=of h1] (g2) {$\times$};
    \node[humanstep, right=of g2] (h2) {Otorisasi\\TTE};
    
    \node[autostep] at (12.5, -4.0) (a3) {Hash-Lock\\Immutable};
    \node[autostep, right=of a3] (a4) {Log Audit\\Trail};
    \node[endstyle_tobe, right=of a4] (e2) {Arsip};

    \draw[arr_tobe] (s2) -- (a1); \draw[arr_tobe] (a1) -- (a2);
    \draw[arr_tobe] (a2.south) -- ++(0,-0.4) -| (h1.north);
    \draw[arr_tobe] (h1) -- (g2);
    \draw[arrno] (g2.north) -- ++(0,0.8) -| node[above, font=\tiny\itshape, text=red] {Batal} (a2.north);
    \draw[arrok] (g2) -- node[above, font=\tiny\itshape] {Setuju} (h2);
    \draw[arr_tobe] (h2.south) -- ++(0,-0.4) -| (a3.north);
    \draw[arr_tobe] (a3) -- (a4); \draw[arr_tobe] (a4) -- (e2);

    \node[timelbl_tobe, below=0.15cm of h1] {SLA: 16 Jam};
    \node[font=\small\bfseries, text=green!50!black] at (9, -6.5) {SIKLUS DIGITAL: < 24 JAM};
    \end{tikzpicture}
    }
\end{frame}

% SLIDE 10: PATTERNS OF INTEGRATION
\begin{frame}{3 Pola Integrasi: "Bumbu Rahasia" Otomasi}
    \small
    \begin{description}
        \item[\textcolor{plnblue}{1. Trigger-Based Auto-Capture}] \hfill \\
        Begitu status di ERP menjadi \texttt{TECO} (Technical Complete), arsip langsung tercipta otomatis.
        \item[\textcolor{plncyan}{2. Parallel Digital Routing}] \hfill \\
        Persetujuan dikirim serentak ke banyak manajer. Tidak menunggu "Map" fisik berpindah.
        \item[\textcolor{plnorange}{3. SLA-Driven Escalation}] \hfill \\
        Jika manajer abai $>16$ jam, sistem otomatis mengeskalasi ke atasan langsung.
    \end{description}
\end{frame}

% SLIDE 11: SOP EVOLUTION
\begin{frame}[fragile]{Evolusi Paradigma SOP PLN}
    \begin{table}
        \centering
        \scriptsize
        \begin{tabularx}{\textwidth}{@{} X X @{}}
            \toprule
            \textbf{SOP Konvensional (Legacy)} & \textbf{SOP Digital Terintegrasi (Future)} \\
            \midrule
            Fokus pada interaksi fisik manusia. & Orkestrasi sistem dan manusia. \\
            Inisiatif manual oleh staf. & \textbf{System-Event Trigger} otomatis. \\
            Validasi tanda tangan basah & Validasi \textbf{Audit Trail} \& \textbf{TTE}. \\
            Risiko: Dokumen tercecer/hilang. & Risiko: Terproteksi dalam \textbf{Immutable Cloud}. \\
            Audit: Tahunan / Sampling. & Audit: Real-time via Dashboard. \\
            \bottomrule
        \end{tabularx}
    \end{table}
    \begin{orangebox}[Insight Pimpinan]
        SOP Digital bukan buku manual sistem, tapi \textbf{Protokol Korporat} untuk mengamankan aset informasi.
    \end{orangebox}
\end{frame}

% SLIDE 12: ANATOMY OF DIGITAL SOP
\begin{frame}{Anatomi SOP Digital Terintegrasi}
    \vspace{0.2cm}
    \begin{columns}
        \column{0.5\textwidth}
        \scriptsize
        \begin{itemize}
            \item \textbf{Scope Interoperabilitas:} Batasan sistem penarik data.
            \item \textbf{Dasar Hukum:} UU ITE \& PP 71/2019.
            \item \textbf{Definisi Digital:} Token, Hash, RBAC, Metadata.
            \item \textbf{System Custodian:} Entitas penanggung jawab mesin.
        \end{itemize}
        \column{0.5\textwidth}
        \scriptsize
        \begin{itemize}
            \item \textbf{System Trigger Points:} Peristiwa sistemik pemicu arsip.
            \item \textbf{The Validation Route (SLA):} Jalur persetujuan \& eskalasi.
            \item \textbf{Exception \& Fallback:} Protokol saat sistem \textit{down}.
            \item \textbf{Immutable Logic:} Jaminan anti-rekayasa.
        \end{itemize}
    \end{columns}
\end{frame}

% SLIDE 13: METADATA EVIDENCE FRAMEWORK
\begin{frame}{14 Field Metadata: Kerangka Bukti Hukum}
    \vspace{0.3cm}
    \begin{columns}
        \column{0.5\textwidth}
        \begin{premiumbox}[Evidence \& Compliance (Otentisitas)]
            \tiny
            \begin{itemize}
                \item \texttt{DOC\_ID}: ID Unik.
                \item \texttt{TIMESTAMP}: Waktu mutlak (ISO 8601).
                \item \texttt{HASH}: Segel digital.
                \item \texttt{ROLE}: Jabatan penandatangan.
                \item \texttt{AUDIT\_LOG\_ID}: Jejak sistem.
            \end{itemize}
        \end{premiumbox}
        
        \column{0.5\textwidth}
        \begin{orangebox}[Discovery \& Retrieve (Nilai Guna)]
            \tiny
            \begin{itemize}
                \item \texttt{ASSET\_ID}: Referensi aset.
                \item \texttt{LOCATION}: Kode GIS.
                \item \texttt{SEARCH\_TAGS}: Kata kunci.
                \item \texttt{RETENTION}: Jadwal JRA.
            \end{itemize}
        \end{orangebox}
    \end{columns}
\end{frame}

% SLIDE 14: LEGAL STANDING
\begin{frame}{Keabsahan Hukum: TTE vs TTD Basah}
    \begin{columns}
        \column{0.6\textwidth}
        \begin{premiumbox}[PP No. 71 Tahun 2019]
            \small
            \textbf{TTE Tersertifikasi} memiliki kekuatan hukum yang setara dengan tanda tangan basah secara absolut.
            \begin{itemize}
                \scriptsize
                \item \textbf{Autentisitas:} Terhubung hanya ke pemilik.
                \item \textbf{Integritas:} Setiap perubahan terdeteksi.
                \item \textbf{Nirkali:} Tidak bisa disangkal.
            \end{itemize}
        \end{premiumbox}
        
        \column{0.4\textwidth}
        \centering
        \stepicon{plngreen}{\checkmark} \\[0.2cm]
        \tiny \textbf{VALID \& SECURE}
    \end{columns}
\end{frame}

% SLIDE 15: KPI & DASHBOARD
\begin{frame}{Dashboard Monitoring: Mengukur Kesuksesan}
    \begin{columns}
        \column{0.25\textwidth}
        \centering
        \stepicon{plnblue}{1} \\ \textbf{Cycle Time} \\ \tiny $\downarrow$ 80\%
        
        \column{0.25\textwidth}
        \centering
        \stepicon{plncyan}{2} \\ \textbf{First-Acc.} \\ \tiny $\uparrow$ 95\%
        
        \column{0.25\textwidth}
        \centering
        \stepicon{plnorange}{3} \\ \textbf{Retrieve} \\ \tiny < 30 Detik
        
        \column{0.25\textwidth}
        \centering
        \stepicon{plngreen}{4} \\ \textbf{Compliance} \\ \tiny 100\%
    \end{columns}
\end{frame}

% SLIDE 16: FINANCIAL ROI
\begin{frame}{Analisis ROI: Investasi vs Manfaat}
    \begin{table}
        \centering
        \scriptsize
        \begin{tabularx}{\textwidth}{@{} X r X r @{}}
            \toprule
            \textbf{Investasi} & \textbf{Biaya} & \textbf{Manfaat} & \textbf{Nilai} \\
            \midrule
            Konfigurasi & Rp 75 Jt & Productivity & Rp 756 Jt \\
            SDM & Rp 36 Jt & Cost Avoidance & Rp 52 Jt \\
            Pilot & Rp 50 Jt & Risk Mitigation & Rp 200 Jt \\
            \midrule
            \textbf{Total} & \textbf{Rp 161 Jt} & \textbf{Total} & \textbf{Rp 1,04 M} \\
            \bottomrule
        \end{tabularx}
    \end{table}
    \centering
    \begin{premiumbox}
        \centering
        \small \textbf{ROI: 550\%} \quad | \quad \textbf{PAYBACK: < 2 BULAN}
    \end{premiumbox}
\end{frame}

% SLIDE 17: ROADMAP & CHANGE MANAGEMENT
\begin{frame}{Rencana Aksi \& Budaya Kerja}
    \begin{columns}
        \column{0.5\textwidth}
        \begin{premiumbox}[30-Day Roadmap]
            \tiny
            \begin{itemize}
                \item \textbf{M1:} API Integration \& Training.
                \item \textbf{M2:} Parallel Run.
                \item \textbf{M3:} Cut-Off Physical Docs.
                \item \textbf{M4:} Standardize SOP.
            \end{itemize}
        \end{premiumbox}
        
        \column{0.5\textwidth}
        \begin{orangebox}[Program "Duta Kearsipan"]
            \small Mendukung \textbf{Budaya Digital} di tiap unit melalui pelatihan dan pendampingan sebaya.
        \end{orangebox}
    \end{columns}
\end{frame}

% SLIDE 18: CLOSING
\begin{frame}
    \vspace{1.5cm}
    \begin{center}
        \huge \textcolor{plnblue}{\textbf{TERIMA KASIH}}\\[0.5cm]
        \large Mari Menuju Ekosistem Informasi yang Adaptif dan Otentik! \\[1.5cm]
        \normalsize \textit{Optimalisasi hari ini untuk Keandalan hari esok.}
    \end{center}
\end{frame}

\end{document}
